\chapter{Detección de Datos Anómalos. Problema de la falta de simetría. Transformaciones.}

\section{Transformaciones}
\noindent
Imaginémonos que hemos detectado anomalías en una distribución, que pueden ser de dos tipos:
\begin{enumerate}
    \item Existencia de datos anómalos.
    \item Existencia de asimetría.
\end{enumerate}
Estas pueden suceder ambas a la vez.\\

\noindent
Si aplicamos una transformaciona a los datos $T(X)$, es posible que corrijamos la asimetría y la existencia de los datos anómalos. Cuando aplicamos una transformación debemos tener en cuenta cómo se afectan las distancias entre los datos.\\

\noindent
La única distancia que no cambia la distancia entre los puntos es una transformación lineal, que solo cambia la escala de la distancia, mediante traslaciones y/o rotaciones. Los tipos de transformaciones que vamos a tener son:
\begin{itemize}
    \item Transformaciones lineales, para cambiar los dominios de las transformaciones.
    \item Transformaciones monótonas no lineales.
    \item Transformaciones no monótonas.
\end{itemize}

\noindent
Sabemos que siempre es posible tener una transformación lineal $h:[a,b]\to [c,d]$. Las transformaciones cóncavas y convexas permiten transformar distancias, por ejemplo para valores pequeños permite acercarlos tras la transformación y para valores grandes permite alejarlos.

% El miércoles pone un ejercicio sin IA
